\documentclass[11pt]{article}
\usepackage[T1]{fontenc}
\usepackage{lmodern}
\usepackage[margin=0.95in]{geometry}
\usepackage{amsmath,amssymb,amsthm,mathtools}
\usepackage{booktabs,longtable,array}
\usepackage{microtype}
\usepackage{xcolor}
\usepackage{enumitem}
\usepackage{hyperref}
\hypersetup{colorlinks=true,linkcolor=blue!45!black,urlcolor=blue!45!black,citecolor=blue!45!black,pdftitle={Polynomial shears and sign recovery for Diophantine consistency},pdfauthor={AI-assisted research prepared for the ProveIt project}}
\newtheorem{theorem}{Theorem}[section]
\newtheorem{lemma}[theorem]{Lemma}
\newtheorem{proposition}[theorem]{Proposition}
\newtheorem{corollary}[theorem]{Corollary}
\theoremstyle{definition}
\newtheorem{definition}[theorem]{Definition}
\theoremstyle{remark}
\newtheorem{remark}[theorem]{Remark}
\newcommand{\Z}{\mathbb Z}
\newcommand{\N}{\mathbb N}
\newcommand{\Pos}{\mathbb Z_{>0}}
\newcommand{\Con}{\operatorname{Con}}
\newcommand{\ZFC}{\mathsf{ZFC}}
\newcommand{\dd}{\operatorname{deg}}
\newcommand{\wt}{\operatorname{wt}}
\newcommand{\file}[1]{\nolinkurl{#1}}
\setlist{itemsep=3pt,topsep=5pt}
\setlength{\emergencystretch}{2em}
\title{Polynomial shears and sign recovery\\for Diophantine consistency}
\author{Research prepared for Vladimir Reshetnikov's ProveIt project\\\small AI-assisted mathematical development and executable audits}
\date{23 October 2026}

\begin{document}
\maketitle
\begin{abstract}
We construct explicit improvements to the currently supplied Diophantine
forms of $\Con(\ZFC)$. A triangular polynomial change of the final integer
coordinate lowers the degree of a ten-integer equation from $1\,245\,324$
to $1\,179\,788$, $1\,114\,252$, or $1\,081\,125$, using respectively
$484$, $487$, and $521$ binary arithmetic operations. The predecessor uses
$480$. The first two substitutions have very short factored formulas; the
third attains a lower bound for this particular family of final-coordinate
translations. Integral changes of coordinates and monic elimination also
give equations in $23,21,20,19,18$ unrestricted integers at degrees
$10,16,20,32,88$. A separate positive-domain argument gives eleven positive
unknowns at degree $211$ in $112$ operations. All these are fixed ordinary
polynomials, with actual coefficient payloads, literal arithmetic data,
proofs, and reproducible checks.

A further section develops natural-number selectors and direct proof-checker
graphs. These improve an alternative proof-certificate source but do not
constitute a completed new global consistency equation. We do not claim a
globally simplest equation, a reduction below nine positive unknowns, or a
new end-to-end proof-assistant certificate. The completed equations inherit
the precisely identified native/compiler/refutation calibration; the new
algebraic reductions and their additional sign argument are proved here.
\end{abstract}
\tableofcontents

\section{Results, comparison, and meaning of simplicity}

The goal is an explicit fixed polynomial $P$ and a specified domain $D$ for
which ordinary mathematics formalizable in ZFC proves
\begin{equation}\label{eq:goal}
 \ZFC\vdash
 \bigl(\neg\exists\boldsymbol x\in D\;P(\boldsymbol x)=0\bigr)
 \longleftrightarrow \Con(\ZFC).
\end{equation}
Consistency here is the syntactic assertion that there is no finite proof
of contradiction in the repository's calibrated recursive ZFC calculus.
No choice of coefficient depends on the truth of consistency.

The reviewed checkout is ProveIt commit
\file{c1a102936aca437e9999530a0111e6ddab56dc6d}. It includes the placed
sources 31--75 and the five archives then in \file{docs/incoming}.
In particular, the comparisons below include the incoming structural,
joint-guard, centering, and radical-descent results; comparing only with
the placed wave-ten README would overstate the improvement.

\begin{table}[ht]
\centering
\small
\begin{tabular}{@{}lrrl@{}}
\toprule
Domain & Unknowns & Exact degree & Additional measurement\\
\midrule
$\Z$ & 10 & $1\,179\,788$ & 484 operations\\
$\Z$ & 10 & $1\,114\,252$ & 487 operations\\
$\Z$ & 10 & $1\,081\,125$ & 521 operations\\
$\Pos$ & 11 & 211 & 112 operations\\
$\Z$ & 23 & 10 & 11 squared residuals\\
$\Z$ & 21 & 16 & 9 squared residuals\\
$\Z$ & 20 & 20 & 8 squared residuals\\
$\Z$ & 19 & 32 & 7 squared residuals\\
$\Z$ & 18 & 88 & 6 squared residuals\\
\bottomrule
\end{tabular}
\caption{Completed new equations in this delivery. Operations mean live
binary $+,-,\times$ gates, including scalar products and squaring; fixed
integer leaves and reuse are free. No optimized operation claim is made
for the five residual-system equations.}
\label{tab:results}
\end{table}

The useful short ten-integer choice is the degree-$1\,114\,252$ equation:
it saves $131\,072$ degrees, about $10.53\%$, for seven extra gates.
The degree-$1\,081\,125$ choice saves $164\,199$ degrees, about $13.19\%$,
at a larger arithmetic cost. The eleven-positive equation improves the
incoming degree-$257$, 114-operation construction in both degree
and operation count. The new $19$-integer equation has degree $32$, compared with the
incoming $19$-integer degree-$88$ equation; the latter also yields a new
$18$-integer degree-$88$ form.

There is no scalar definition of ``simplest'' inherent in the problem.
Number of unknowns, domain, total degree, circuit size, coefficient size,
and verification burden are distinct. In particular, an equation in
$\Pos^{11}$ is not an eleven-unrestricted-integer equation. The established
nine-positive and ten-integer variable counts remain relevant. Sun's
ten-integer theorem is a recent primary reference for the latter type of
compression \cite{sun2026}; Jones's paper is the classical reference for
the nine-positive route \cite{jones1982}. Our new changes do not improve
those general variable bounds.

\section{A general polynomial-shear principle}

\begin{lemma}[Triangular shear]\label{lem:shear}
Let $P\in\Z[\boldsymbol v,m]$ and $R\in\Z[\boldsymbol v]$. Set
\[
 P_R(\boldsymbol v,m)=P(\boldsymbol v,m+R(\boldsymbol v)).
\]
Then $(\boldsymbol v,m)\mapsto(\boldsymbol v,m+R(\boldsymbol v))$ is a
bijection of $\Z^{n+1}$ with inverse
$(\boldsymbol v,m)\mapsto(\boldsymbol v,m-R(\boldsymbol v))$.
It restricts to a bijection between the integer zeros of $P_R$ and of $P$.
\end{lemma}
\begin{proof}
Both maps are integral polynomial maps, and their two compositions are the
identity. Substituting into the displayed definition of $P_R$ proves the
zero-set statement. No positivity, divisibility, or nonvanishing premise
is needed.
\end{proof}

The following norm shape makes a suitable shear especially effective:
\begin{equation}\label{eq:norm}
 P=\bigl((\Sigma m-\Theta)^2-\Sigma^2(a_0+a_1)\bigr)^2
       -4\Sigma^4a_0a_1.
\end{equation}
Every displayed symbol other than $m$ is a polynomial in the other
coordinates. Suppose $R$ is a polynomial approximation to one radical in
the purely algebraic sense that $\varepsilon=a_1-R^2$ has smaller degree
than $a_1$. Writing $Z_0=\Sigma m-\Theta$, the shifted equation is exactly
\begin{equation}\label{eq:centered-norm}
 P_R=\bigl(Z_0^2+2\Sigma RZ_0-\Sigma^2(a_0+\varepsilon)\bigr)^2
       -4\Sigma^4a_0a_1.
\end{equation}
This follows by expanding $(Z_0+\Sigma R)^2$ and canceling $R^2-a_1$.
An approximation inequality for square roots is unnecessary. The identity
holds on every signed integer assignment, including assignments where
$a_1$ is negative or zero.

\begin{proposition}[Degree accounting]\label{prop:degree}
Let $s,z,r,d_0,d_1,e$ be upper bounds on the degrees of
$\Sigma,Z_0,R,a_0,a_1,\varepsilon$, respectively. Then
\[
 \dd P_R\leq
 \max\{2\max(2z,s+r+z,2s+d_0,2s+e),\ 4s+d_0+d_1\}.
\]
When one of these contributions strictly dominates all other contributions
and its leading homogeneous part is nonzero, the bound is attained.
\end{proposition}
\begin{proof}
Apply $\dd(fg)\leq\dd f+\dd g$ and
$\dd(f+g)\leq\max(\dd f,\dd g)$ to \eqref{eq:centered-norm}.
Under strict dominance, lower-degree summands cannot cancel the highest
homogeneous part.
\end{proof}

\section{The fixed ten-integer source and final norm}

\subsection{Exact data rather than an unspecified universal parameter}

The predecessor is the incoming Round11 H1868 norm-saturated equation
\cite{round11}. Its literal 480-operation graph is included unchanged as
\file{norm10/upstream/integer10.json}. Its proof, source-packing proof,
four native integer payloads, twelve actual live fixed integer leaves, and
licenses are retained. The new graphs are complete ordinary polynomial
graphs; the predecessor file is used to construct and independently check
them, rather than being an existential parameter in the new equations.

The ten unrestricted integer variables are
\[
 t,g_0,c_0,k_0,\ell,w,x,y,z,m.
\]
Use the following source abbreviations:
\[
 b=6t+4,\qquad g=bg_0,\qquad k=b(3k_0+2),\qquad
 B_c=2^{2134940}b^2.
\]
The finite weight list, target list, and three fixed sparse kernel
polynomials are the literal data in
\file{norm10/data/code.json} and \file{norm10/data/kernels.json}.
Write them as $w_i,q_j,D_0,D_1,D_2$, respectively, and set
\begin{align*}
 e_c&=\sum_i B_c^{w_i},& U_c&=e_c+g,\\
 G_c&=D_2(B_c)U_c^2+D_1(B_c)U_c+D_0(B_c),&
 T_c&=\sum_j B_c^{q_j},\\
 N_0&=B_c^{117},&N&=B_c^{1986},\\
 M_c&=N_0-1-(b-1)e_c,\\
 S_c&=g+N_0(2G_c+T_c),&
 T'_c&=M_c+N_0(B_c-2)T_c-2,\\
 X&=N(S_c+T'_c+1)-S_c-1,\\
 Y&=2^{34159017}b B_c^{2109},&\mathcal C&=bB_c^{116}.
\end{align*}
These are fixed finite sums and fixed nonnegative powers. Kernel
coefficients are explicitly materialized integers, determined by the
whole inherited native tuple. Neither a program index chosen by a
nonconstructive existence theorem nor a variable exponent occurs in the
equation.

\subsection{Polynomial registers}

Define
\begin{align*}
 &L=\ell Y,\quad U=2LX,\quad V=4gwY,\quad W=bw,\\
 &h_A=LX(V+1),\quad A=2h_A,\quad B_1=2X+1,\quad C=6c_0-1,\\
 &P_h=(LX)UV,\quad K=X+1+k(2P_h-2),\\
 &S=2A-5,\quad T=3WC-2(W^2-1),\\
 &Q=(P_h^2-1)K^2+1,\\
 &\Delta=8\mathcal C^3gK^2
       -g^2\bigl(32(C-KL)^2\mathcal C^3+g^2K^2\bigr),\\
 &\xi=zx,\quad D=(A^2-4)C^2+4,\\
 &E_0=C^2\xi,\quad E=DE_0,\quad E_2=E_0E,\\
 &F=4(A^2-4)E_2+1,\quad G=(h_A-1)(F-1)-1,\\
 &\Omega=16E,\quad H=C+F(B_1-C+\Omega y),\\
 &I=(G^2-1)H^2+1,\\
 &\rho=T+Sy,\quad \Sigma=\xi S^5,\\
 &\Theta=(\rho^3-\rho S^2-S^3)(\rho^2+23S^2),\\
 &r=bQ,\quad a_0=r+(9\Delta-3)(rz)^2,\quad a_1=DFI.
\end{align*}
The half-parameter $h_A$ is explicitly the integer polynomial above.
No division is being used to define it. Every register is an abbreviation;
only the original ten coordinates are quantified.

The supplied predecessor theorem says that the polynomial \eqref{eq:norm}
has an integer zero exactly when the calibrated calculus has a finite
refutation. We inherit that theorem, including its Pell and native-source
proof obligations. The present section's transformations do not change
any of its source parameters or coefficient payloads.

\section{Three new ten-integer equations}

\subsection{Four additional gates}

Take
\begin{equation}\label{eq:R0}
 R_0=2AEGH.
\end{equation}
The new equation is \eqref{eq:norm} with $m$ replaced everywhere by
$m+R_0$. The essential cancellation is
\begin{equation}\label{eq:first-defect}
 DF-(2AE)^2=D-16E^2.
\end{equation}
Indeed $DF=4(A^2-4)E^2+D$. Consequently,
\begin{equation}\label{eq:R0def}
 a_1-R_0^2=(D-16E^2)G^2H^2+DF(1-H^2).
\end{equation}
The shared register $2A$ already occurs in $S=2A-5$. Only three new
multiplications and the addition to $m$ are needed. The exact graph has
$484$ live operations.

\subsection{Seven additional gates}

A better short substitution is
\begin{equation}\label{eq:R1}
 R_1=2AE_0(D-2C^2)GH.
\end{equation}
Put $P_1=2AE_0(D-2C^2)$ and $u=C^2$. Direct polynomial expansion gives
\begin{equation}\label{eq:second-defect}
 DF-P_1^2
 =D+16E_0^2\bigl(u(3u-4)A^2-16(u-1)^2\bigr).
\end{equation}
For a transparent derivation, put $A=2h_A$ and
$D=4D_*\!,\ D_*=u(h_A^2-1)+1$. Then
\[
 DF=256(h_A^2-1)u^2D_*^2\xi^2+4D_*,\qquad
 P_1=8u\xi h_A\bigl(2uh_A^2+2-3u\bigr).
\]
The $h_A^6$ and $h_A^4$ terms cancel on subtraction; the remaining
expression is exactly \eqref{eq:second-defect}. Hence
\begin{equation}\label{eq:R1def}
 a_1-R_1^2=(DF-P_1^2)G^2H^2+DF(1-H^2).
\end{equation}
Compared with \eqref{eq:R0def}, two further powers of the high-degree
parameter $A$ disappear. The resulting graph has $487$ operations.

\subsection{A six-term polynomial part}

Let $h=h_A$, $u=C^2$, and $J=B_1-C$. Define the following short polynomial:
\begin{align}
 T_*={}&\xi J\bigl(u(2h^2-11)+6\bigr)\nonumber\\
 &+16u\xi^2y\bigl(u^2(8h^4-52h^2+143)
                 +u(32h^2-176)+48\bigr)+3y,\label{eq:Tstar}\\
 R_6={}&32768u^7\xi^4h^9\bigl((h-1)T_*-2y\bigr).
 \label{eq:R6}
\end{align}
This polynomial has integer coefficients. It is obtained by matching the
six highest powers $h^{14},\ldots,h^9$ in a formal square root of $a_1$.
That discovery description is not needed as a correctness assumption:
the finite formula \eqref{eq:R6} and exact expansion of $a_1-R_6^2$ are
the certificate.

For reproducibility, the file \file{norm10/ROOT_CERTIFICATES.json}
gives the full coefficient list of $R_6$ and of its defect. In the five
independent indeterminates $(h,B_1,C,\xi,y)$, $R_6$ has $26$ monomials and
the defect has $2306$. The independent standard-library checker multiplies
the polynomials and checks every coefficient. The defect has degree at
most $22$ in $h$; its relevant total-degree bound after source substitution
is proved next. The complete literal graph uses $521$ live operations.

\subsection{Exact total degrees}

The source has
\[
 \dd h_A=\dd A=16384,\qquad \dd B_1=7942,
 \qquad \dd C=1,\quad\dd\xi=2,\quad\dd y=1.
\]
The weighted degrees in these five indeterminates therefore give valid
upper bounds after substitution. For the actual source the relevant
degrees are
\begin{center}
\begin{tabular}{@{}lrlr@{}}
\toprule
Register&Degree&Register&Degree\\
\midrule
$D$&32770&$E$&32774\\
$F$&65546&$G$&81930\\
$H$&98321&$a_1$&458818\\
$\Sigma$&81922&$\Theta$&81925\\
$a_0$&294619&$\Sigma m-\Theta$&81925\\
\bottomrule
\end{tabular}
\end{center}
Every proposed $R$ has degree $229409$. The cancellation identities give
\[
 \dd(a_1-R_0^2)=426050,\qquad
 \dd(a_1-R_1^2)=393282,
\]
and the full finite defect certificate gives
\[
 \dd(a_1-R_6^2)\leq368449.
\]
Applying Proposition~\ref{prop:degree} yields respectively
\[
 1\,179\,788,\qquad1\,114\,252,\qquad1\,081\,125.
\]
These are exact degrees, not merely syntactic upper bounds.
For $R_0$ and $R_1$, the highest term comes from the square of
$-\Sigma^2(a_1-R^2)$. For $R_6$, the square in
\eqref{eq:centered-norm} has degree at most $1\,064\,586$, whereas
\[
 \dd(-4\Sigma^4a_0a_1)
 =4(81922)+294619+458818=1\,081\,125.
\]
There can be no cancellation between these unequal degrees.

The checker evaluates the claimed leading homogeneous coefficient after
the diagonal specialization of all ten variables at a common parameter,
using the actual fixed integers reduced modulo three explicitly supplied
primes. For $R_6$, the residues at $1000003,1000033,1000037$ are
$347224,762141,563833$, respectively. Each nonzero residue proves that the
integer leading coefficient is nonzero. A single nonzero residue suffices;
the additional residues are transcription controls. The leading terms for
the first two variants are likewise certified. This proves attainment of
the upper bounds without constructing million-degree expanded polynomials.

\begin{theorem}\label{thm:ten}
The three literal polynomials in
\file{norm10/integer10_linear.json},
\file{norm10/integer10_cubic.json}, and
\file{norm10/integer10_sixterm.json} have the variable counts, exact degrees,
and operation counts shown in Table~\ref{tab:results}. Each is related to
the pinned predecessor by an integral triangular polynomial automorphism.
Composing with the inherited predecessor theorem gives \eqref{eq:goal}
for each of these polynomials on $\Z^{10}$.
\end{theorem}
\begin{proof}
The literal graphs implement \eqref{eq:norm} with the respective formulas
\eqref{eq:R0}, \eqref{eq:R1}, and \eqref{eq:R6} substituted for the shift.
Lemma~\ref{lem:shear} proves both existence implications, on all signed
assignments. The preceding degree calculation and finite graph counts prove
the quantitative statements. The source-to-refutation equivalence is
exactly the inherited theorem, and negating existence yields the assertion.
\end{proof}

\subsection{A lower bound for this restricted transformation family}

There is a useful stopping criterion. Keep the entire source, $\Sigma$,
$\Theta$, $a_0$, and $a_1$ fixed, and allow only polynomial translations of
$m$ with integer coefficients. In every resulting equation,
\[
 P_R=\mathcal B_R^2-4\Sigma^4a_0a_1
\]
for some integer polynomial $\mathcal B_R$. The second summand has the
odd degree $1\,081\,125$. The nonzero square $\mathcal B_R^2$ has even
degree. If its degree is larger, it cannot cancel against the second
summand; if it is smaller, it cannot cancel that summand either. Equality
of degrees is impossible by parity. The case $\mathcal B_R=0$ is immediate.
Thus every such translation has degree at least $1\,081\,125$.
The six-term choice attains this bound.

This is an exact optimality statement for a sharply specified transformation
family. It is not a lower bound for arbitrary polynomial coordinate
changes, alternative norm representations, source encodings, or all
Diophantine equations expressing consistency.

\section{Fewer unrestricted integer variables}
\label{sec:integer}
\subsection{Double centering and a transport pivot over the integers}
\label{sec:integer-projections}

The following five equations use unrestricted integer variables and the
same six fixed coefficients throughout.  Each is a sum of squares of the
residuals in its accompanying JSON packet.
\begin{center}
\begin{tabular}{lrrr}
\hline
Packet & Variables & Residuals & Exact degree\\
\hline
\texttt{centered23\_degree10} &23&11&10\\
\texttt{pivot21\_degree16} &21&9&16\\
\texttt{pivot20\_degree20} &20&8&20\\
\texttt{pivot19\_degree32} &19&7&32\\
\texttt{pivot18\_degree88} &18&6&88\\
\hline
\end{tabular}
\end{center}
The inspected incoming Joint Guards package had the corresponding points
$23/12$, $21/20$, $20/32$, and $19/88$. The incoming Radical Guard Descent
$22/14$ form remains another useful tradeoff. These improvements do not
lower the ten-integer high-degree variable bound.

All new coordinate transformations below are integer-polynomial bijections
of the indicated parent zero sets. Consequently they preserve the parents'
ordinary, ZFC-formalizable equivalence with finite syntactic ZFC refutability.
Their unsolvability is therefore ZFC-provably equivalent to
$\operatorname{Con}(\mathrm{ZFC})$. This is an ordinary mathematical
composition; no new end-to-end kernel certificate or object-language ZFC
derivation is claimed.

The constituent two-root centering appears in incoming Budget and Centering,
and the monic transport pivot appears in incoming Round11. The new result is
their application to the newer Joint Guards systems. No historical priority
or global optimality claim is made.

\subsection{Fixed coefficients}
The coherent tuple is
\[
 \beta=2^{1064921}-1,\quad K_0,M_C,M_F,\quad
 A_*=2^{1113},\quad B_*=2^{666104}.
\]
Here $K_0,M_C,M_F$ are exactly the source-34
\texttt{mixed-parity-unrounded-1064921} native tuple, with $M_F$ already
shifted once. The six actual minimal unsigned little-endian byte strings
are supplied, compressed with gzip, with a complete binding manifest. Their
combined raw size is $615{,}883$ bytes. These are fixed numeral leaves,
not quantified parameters.

\subsection{The 23-variable equation}
Take independent integer coordinates
\[
 (J,F,\alpha,z,f,h,r,v,w,\tau,\eta,\zeta,\nu,\gamma,o,j,E,a,I,t,p_1,p_2,p_3).
\]
Define polynomial abbreviations only:
\begin{align*}
 q&=\beta(3J+1)+1,& W&=4(3w+1),& X&=qW,\\
 Y&=a-E,& k&=\eta+\zeta,& c&=Yk+\eta,\\
 V&=fo-c,& C&=q-F-\alpha,& H&=4a+3,\\
 \Delta&=a^2+H,& T&=E(c-\eta),& S&=X+\gamma H.
\end{align*}
Also put
\begin{align*}
 R&=((q-1)(q-F)+\alpha)(q^2-1)+(M_C+qM_F)(3J+1),\\
 G&=(A_*F)(B_*C)-(q-A_*F-B_*C)^2,\\
 Z&=\eta\zeta,\qquad
 \mathrm{Pos}=1+p_1^2+p_2^2+p_3(p_3+1).
\end{align*}
The eleven residuals are
\begin{align*}
 B_1&=a-E-(2v+1)q^3,\\
 B_2&=E-XY,\\
 B_3&=\tau^2+(2\tau-k)T-1,\\
 B_4&=2acS+S^2-Hc^2-1,\\
 B_5&=k-1-hE-R,\\
 B_6&=(K_0+W)C-F-z(q-1),\\
 B_7&=I-(2r+1)c^2,\\
 B_8&=I^2-\Delta(f^2-1),\\
 B_9&=(V+\nu)^2-I^2\nu(2V+\nu)-1,\\
 B_{10}&=V+k-1-cj-hE,\\
 B_{11}&=GZ-(t^2-G-Z+1)^2-\mathrm{Pos}.
\end{align*}
Thus the complete ordinary polynomial equation is
\[
 \boxed{\displaystyle P_{23}=\sum_{i=1}^{11}B_i^2=0.}
\]
Every displayed abbreviation is inlined in the supplied packet; it
introduces neither a variable nor a side condition.

For the proof, the incoming 25-variable octic has independent $c,V,y$ and
uses $k_0=1-V+jc+hE$ in its $c$ equation and first norm. It retains the
row $L=k-k_0$. Replace $k_0$ by $k$ using polynomial multiples of this
retained row. Then make the triangular root translations
\[
 \tau_{\mathrm{old}}=\tau+E(c-\eta),\qquad
 y_{\mathrm{old}}=V+\nu,
\]
and eliminate the two monic definitions $c=(a-E)k+\eta$ and $V=fo-c$.
The three cancellation identities are
\begin{align*}
 (\tau+T)^2-T(k+T)-1
   &=\tau^2+(2\tau-k)T-1,\\
 (ac+S)^2-(a^2+H)c^2-1
   &=2acS+S^2-Hc^2-1,\\
 I^2\bigl(V^2-(V+\nu)^2\bigr)+(V+\nu)^2-1
   &=(V+\nu)^2-I^2\nu(2V+\nu)-1.
\end{align*}
In the literal parent order, its thirteen restored rows are exactly
\[
 B_1,B_2,(a-E)B_{10},B_3+TB_{10},B_4,B_5-B_{10},
 B_6,0,B_7,B_8,B_9,B_{10},B_{11}.
\]
This gives both directions of the zero-set equivalence without any
division. The inverse root translations and coordinate projection are
unique. The residual degrees are
\[
 4,3,4,5,4,2,5,4,5,3,4.
\]
Only $B_7$ contains $r$, and
\[
 [r^2a^4\eta^4]P_{23}=4.
\]
Hence the total degree is exactly ten on the actual fixed-coefficient slice.

\subsection{A transport pivot that costs no dominant degree}
All four remaining parents contain
\[
 (K_0+W)(q-F-\alpha)-F-z(q-1)=0.
\]
Change the complement coordinate by $C=q-F-\alpha$, then solve the
now monic row for $F$:
\[
 F=(K_0+W)C-z(q-1),\qquad \alpha=q-F-C.
\]
Delete that row and $F$, and substitute everywhere else. Both the map and
its inverse use only integer polynomials, so this is a bijection of integer
zero sets. The packing register becomes
\[
 R=(q-F)(q^3-q)-C(q^2-1)+(M_C+qM_F)(3J+1),
\]
of degree at most five. The determinant $G$ has degree four, and its
contribution to the joint-sign residual has degree at most eight.
The larger Pell/sign residuals therefore already pay the resulting degree
cost in the selected parents. This yields the $21/16$, $20/20$, $19/32$,
and $18/88$ packets in the table.

\subsection{The complete 18-variable endpoint}
Take independent integer coordinates
\[
 (J,C,z,f,h,r,v,w,\tau,y,\gamma,o,j,c,t,p_1,p_2,p_3).
\]
In addition to $q,W,X$ and the transport pivot above, define
\begin{align*}
 Y&=(2v+1)q^3,& E&=XY,& a&=(X+1)Y,\\
 V&=fo-c,& k&=1-V+jc+hE,& \eta&=c-Yk,\\
 \zeta&=k-\eta,& H&=4a+3,& \Delta&=a^2+H,\\
 I&=(2r+1)c^2,& T&=E(c-\eta),& D&=X+ac+\gamma H.
\end{align*}
Use the same $R,G,Z,\mathrm{Pos}$ definitions. The six residuals are
\begin{align*}
 A_1&=\tau^2-T(k+T)-1,\\
 A_2&=D^2-\Delta c^2-1,\\
 A_3&=cj-V-R,\\
 A_4&=I^2-\Delta(f^2-1),\\
 A_5&=I^2(V^2-y^2)+y^2-1,\\
 A_6&=GZ-(t^2-G-Z+1)^2-\mathrm{Pos}.
\end{align*}
The equation is
\[
 \boxed{\displaystyle P_{18}=\sum_{i=1}^{6}A_i^2=0.}
\]
These six residuals are completely inlined in its packet, so the count
18 includes all quantified variables. Retaining $\eta$ as a nineteenth
independent coordinate and adding $c-\eta-Yk=0$ gives the degree-32 sibling.

\subsection{Degree certificates and verification}
The four pivot equations have the following exact nonzero coefficients:
\begin{align*}
 [a^4f^8o^4]P_{21}&=1,\\
 [f^4v^4J^{12}]P_{20}&=2^4 3^{12}\beta^{12},\\
 [\eta^4h^4w^4v^4J^{16}]P_{19}&=2^{12}3^{20}\beta^{16},\\
 [h^8w^8v^{16}J^{56}]P_{18}&=2^{32}3^{64}\beta^{56}.
\end{align*}
For the last two, the leading forms are
\[
 E_{\mathrm{top}}=1944\beta^4J^4wv,
 \qquad \eta_{\mathrm{top}}\big|_{P_{18}}
   =-104976\beta^7hJ^7wv^2.
\]
In $P_{19}$, the dominant joint-sign part is
$-(\eta hE_{\mathrm{top}})^2$ before the outer SOS square. In $P_{18}$,
it is $-\eta_{\mathrm{top}}^4$. These are the unique contributions at
the respective highest residual degrees. The first two certificates come
from unchanged reduced-auxiliary and strong-Pell rows; the pivoted guard
cannot contribute their selected monomials. Full derivations are given in
the companion proof.

The generator authenticates declaration closure, exact degrees and all
leading coefficients. A separate checker imports no builder: it checks
all thirteen centering identities, every literal pivot identity, the actual
six coefficient byte strings and additional signed modular transcription
diagnostics. The finite tests certify only the algebra and the data stated;
the universal native/Pell/refutation calibration remains an explicitly
inherited ordinary theorem.

The next concrete targets are a positivity certificate that avoids squaring
the degree-eight sign register in $P_{19}$, and an orientation theorem that
avoids the degree-eleven $\eta$ substitution in $P_{18}$. Either would
improve these equations more than further cosmetic factorization.


\section{An eleven-positive equation of degree 211}
\label{sec:positive}
\subsection{Eleven positive unknowns: an auxiliary offset and two root centers lower degree to 211}
\label{sec:positive-eleven-offset}

The following construction has eleven strictly positive integer unknowns,
exact degree $211$, and $112$ live binary arithmetic operations. The
centered predecessor in the incoming joint-guards package has degree $257$
in $114$ operations; the uncentered Pell-descent version has degree
$277$ in $111$ operations. The improvement is therefore a strict improvement in both degree and evaluator size over the centered predecessor.

The new step is the affine substitution
\[
 T_{\rm old}=T+F+\alpha.
\]
It makes the auxiliary quantity positive at an earlier stage of the proof,
so an expensive exceptional-sign term can be removed from the transport
weight. A second center absorbs the strong root, leaving a degree-two
gap bound. This is an existence-preserving substitution on canonical source
witnesses, not a bijection of the entire positive zero sets.

\subsection{The complete polynomial}

Use the coherent compact44 fixed tuple $(\beta,K_0,M_C,M_F)$, where
$\beta=2^{590069}-1$ and $M_F$ is already shifted. The tuple is defined by
the retained finite recipe and actual integer payloads. The unknowns are
\[
 j,F,\alpha,i,T,X,Y,n,\eta,\zeta,\gamma>0.
\]
Every remaining symbol is a polynomial abbreviation. Define
\begin{align*}
 J&=3j+1,\qquad q=\beta J+1,\qquad C=q-F-\alpha,\\
 E&=XY,\qquad k=\eta+\zeta,\qquad c=kY+\eta,\\
 a&=E+Y,\qquad H=4a+3,\qquad\Delta=a^2+H=(a+2)^2-1,\\
 \ell&=EkY,\qquad I=ic^2,\qquad M=q^3(q-1),\\
 R&=((q-1)(q-F)+\alpha)(q^2-1)+(M_C+qM_F)J,\\
 B_t&=(qK_0+X)C-qF,\qquad D_\tau=\ell(\ell+k)+1,\\
 D_f&=1+\Delta I^2,\qquad Q=\Delta^2I^2,\qquad\delta=Q-1,\\
 T_{\rm old}&=T+F+\alpha,\qquad C_f=(a+2)I,\\
 V&=D_f(cT_{\rm old}-R)-c,\qquad L=kT_{\rm old},\\
 W&=1+(q-1)L,\\
 \overline Z&=M(\ell+k+n+L+V+C_f)+(q-1)Y-q^3WB_t,\\
 r_0&=k-R-1,\qquad S=4qY,\qquad U=qS,\\
 D_*&=U(X+ac)+H(U\gamma+SX-r_0),\\
 A_1&=D_*^2-U^2(1+\Delta c^2).
\end{align*}
The three-root norm is
\begin{align*}
 A&=M^2D_\tau,\qquad B_f=M^2D_f,\qquad C_{\rm num}=M^2(QV^2-1),\\
 T_N&=\delta^2AB_f,\qquad R_N=\delta C_{\rm num},\\
 J_N&=\delta(\overline Z^2-A-B_f)-C_{\rm num},\\
 N_c&=J_N^2-4[T_N+R_N(A+B_f)],\\
 N_8&=N_c^2-64T_NR_N\overline Z^2.
\end{align*}
The proposed single polynomial equation is
\begin{equation}\label{eq:offset-p211}
 \boxed{P_{11}=N_8+A_1+N_8A_1=0.}
\end{equation}
Its literal graph is \texttt{p11\_positive\_degree211.json}.

\subsection{Recovering radicals before any transport sign}

Put, only in this proof,
\[
 t=\sqrt{D_\tau},\qquad f=\sqrt{D_f},\qquad
 y=\sqrt{\frac{QV^2-1}{Q-1}}.
\]
The unconditional three-root separation argument of source70 continues to
apply: it requires only $T_{\rm old}>0$ and $I\ge c^2$, not positivity of
$C$. We recall its substance. Let $b=cT_{\rm old}-R$. If $b=0$, then
$V=-c$, $y<c+1$, and $f-t>(a+1)c(c-1)>c+1$. If $b\ne0$, then
\[
 y>|V|\ge D_f-c>2f>f+t.
\]
Here $f>3c^2$ and $f>t$ follow directly from $\Delta\ge15$ and
$c>kY$. Thus one of the three roots strictly dominates the other two.
Their positive rational radicands are pairwise distinct: $D_\tau<D_f<Q$,
and equality of the third radicand to either integer $D\in\{D_\tau,D_f\}$
would imply $Q(D-V^2)=D-1$, impossible since $0<D-1<Q$.

A rational signed sum of three such roots forces them all rational. For
signed roots $x_1,x_2,x_3$ with rational nonzero sum $s$ and squares $a,b,c$,
set
\[
 u=\frac{s^2-a-b-c}{2},\qquad
 v=\frac{u^2-ab-ac-bc}{2s}.
\]
Then
\[
 x_1=\frac{a(a+u)}{as+v},\qquad
 as+v=x_1(x_1+x_2)(x_1+x_3)\ne0,
\]
and the cyclic formulas prove the assertion.

The universal norm identity
\[
 \bigl[(x^2-a-b-c)^2-4(ab+ac+bc)\bigr]^2-64abcx^2
 =\prod_{\varepsilon\in\{-1,1\}^3}
 (x-\varepsilon_1\sqrt a-\varepsilon_2\sqrt b-\varepsilon_3\sqrt c)
\]
therefore shows that $N_8=0$ supplies rational signed roots with
\[
 \frac{\overline Z}{M}=\tau_*+f_*+y_*,\qquad
 \tau_*^2=D_\tau,\quad f_*^2=D_f,\quad
 QV^2-(Q-1)y_*^2=1.
\]
The first two roots are integers. The adjacent-coefficient Pell descent
of source69 makes $y_*$ an integer as well. In that descent, writing
$y_*=a/m$ in lowest terms gives $m^2\mid Q-1$. After replacing $V$ and
$y_*$ by their absolute values at each step, the transformation
\[
 (V,y_*)\longmapsto
 ((2Q-1)V-2(Q-1)y_*,\ (2Q-1)y_*-2QV)
\]
preserves the norm and the denominator, while strictly reducing $|V|>1$.
Strong induction reaches $|V|=1$, where $|y_*|=1$.

Since $M\mid\overline Z$, reduction modulo $q^3$ and $q-1$ now gives
\[
 q^3\mid Y,\qquad q-1\mid B_t.
\]
No sign of $B_t$ has been used. The equation $A_1=0$ also gives an
integral main root $D_*/U$, whose absolute value can be classified.
Writing $\chi_A,\psi_A$ for the usual Pell sequences, positive
classification supplies indices $m_1,p\ge1$ such that
\[
 |\tau_*|=\chi_{2XY^2+1}(m_1),\qquad
 k=2\psi_{2XY^2+1}(m_1),\qquad c=\psi_{a+2}(p).
\]
Because $c\ge3$, $p\ge2$. The case $m_1=1$ would give
$k=2$, $\eta=\zeta=1$, $c=2Y+1$, contradicting
$c\ge2(a+2)\ge4Y+4$. Hence
\begin{equation}\label{eq:offset-early-c}
 c>kY\ge4(2XY^2+1)Y>8XY^3\ge8q^9.
\end{equation}
This is established before $C>0$, before the source bounds on $R$, and
before interpreting any recovered transport gap.

\subsection{The affine offset eliminates the exceptional-sign case}

Define the proof-only bound $\widehat R=q^3(q+F+\alpha)$.
The fixed-mask estimates give unconditionally
\[
 |R|<\widehat R.
\]
Put $h=F+\alpha\ge2$. Since $q\ge16$,
\[
 q^3(q+h)\le q^3(q/2+1)h<q^4h<ch
\]
by \eqref{eq:offset-early-c}. Thus
\[
 cT_{\rm old}-R=cT+ch-R>cT\ge c,
 \qquad V>D_fc-c>1.
\]
The new auxiliary coordinate is therefore positive at the required point
of the proof. This is what permits deletion of the predecessor's
$4qD_f$ exceptional-sign weight.

The same bound $ch>\widehat R$ also gives
\[
 0<cT_{\rm old}-R<c(T+2h),\qquad 0<V<D_fc(T+2h).
\]
Moreover,
\[
 \delta-(\Delta-1)D_f=\Delta(I^2-1)>0,
 \qquad c<k(Y+1),\qquad \Delta-1>2(Y+1).
\]
The last inequality follows from $a\ge2Y$ and $Y\ge1$. Therefore
\[
 0<y-V=\frac{V^2-1}{\delta(y+V)}<\frac{V}{\delta}
 <\frac{c(T+2h)}{\Delta-1}<\frac{k(T+2h)}2<k(T+h)=L.
\]
Also $f<C_f=(a+2)I$, since $C_f^2-D_f=I^2-1>0$. We conclude
\begin{equation}\label{eq:offset-small-L}
 L+C_f+V-f-y>0.
\end{equation}

Let $s=Y/q^3>0$ and $b=B_t/(q-1)\in\mathbb Z$. The recovered norm
identity reads
\[
 \tau_*+f_*+y_*=\ell+k+n+L+V+C_f+s-Wb.
\]
If $b\le0$, its right side is strictly larger than
$|\tau_*|+|f_*|+|y_*|$, because $|\tau_*|<\ell+k$ and
\eqref{eq:offset-small-L} holds. Hence $B_t>0$, which forces $C>0$.
Only now do we invoke the native source bounds $3q<R<q^4-2$.

Set
\[
 f_{\rm old}=|f_*|,\quad i_{\rm old}=i,\quad u=|y_*|-V,\quad
 r=n+L+V+C_f-f_*-y_*.
\]
All are positive integers, with $r>n$. They satisfy the weighted source70
carrier with weight $h_*=L$, its normalized strong norm, its auxiliary
norm, and its adaptive main norm. That inherited interface is sound for
every positive integer $h_*$. Its native/compiler and syntactic-refutation
calibration therefore supplies a finite ZFC refutation.

\subsection{Completeness and the exact cost of the offset}

A finite refutation gives the canonical compact source, with
$R\equiv3\pmod4$, $R>3q$, $F+\alpha<q$, and the earlier bound on $c$.
Choose its normalized strong pair $f,I$, $I=ic^2$, and write
$M_*=\Delta I$. The inherited positive auxiliary reconstruction uses
\[
 V_* =\frac{\chi_{M_*}(R)}{M_*},\qquad
 y_* =\psi_{M_*}(R),\qquad
 T_{\rm old}=\frac{V_*+c+Rf^2}{cf^2}.
\]
It proves the last quantity integral and $y_*>V_*>0$. The stronger bound
needed for the new offset is elementary. Since $R\ge7$,
\[
 V_*\ge M_*^{R-1}\ge M_*^6,\qquad f^2<M_*^2,\qquad M_*\ge c^2,
\]
and hence
\[
 T_{\rm old}>\frac{V_*}{cf^2}>\frac{M_*^4}{c}\ge c^7>q>F+\alpha.
\]
Thus $T=T_{\rm old}-F-\alpha>0$ is integral. No new cofinality assumption
or congruence on the strong Pell pair has been introduced.

The same growth also proves $T_{\rm old}>C_f$: because $I\ge c^2$,
\[
 T_{\rm old}>\frac{M_*^4}{c}=\Delta^4I\frac{I^3}{c}>(a+2)I=C_f.
\]
After these choices compute $L,W$. We have $L=kT_{\rm old}>C_f$ and
$L>k$. The canonical transport has $t_0=B_t/(q-1)\ge q$, $s<k$, and
$\tau>\ell$, so
\[
 r=Wt_0+\tau-\ell-k-s>qW-2k>2L.
\]
Consequently
\[
 n=r-L-V_*-C_f+f+y_*=r-L-C_f+f+(y_*-V_*)>0.
\]
This makes $\overline Z=M(\tau+f+y_*)$ and $N_8=0$; the inherited
adaptive carrier supplies positive $\gamma$ and $A_1=0$. Completeness
is proved on canonical witnesses, which is exactly what an existence
equivalence requires.

Both residuals are squares modulo four. Thus the factorization
\[
 1+P_{11}=(1+N_8)(1+A_1)
\]
shows $P_{11}=0$ if and only if $N_8=A_1=0$: a negative unit would
require a residual $-2$, impossible modulo four.

\subsection{Exact degree and verification}

The degree audit uses an exact universal double-center identity. For
independent variables $z,a,b,c,e,f$, set
\[
 x=z+b+c,\quad B=b^2-e,\quad C=c^2-f,\quad
 N_c^\circ=(x^2-a-B-C)^2-4(aB+aC+BC),
\]
\[
 N^\circ=(N_c^\circ)^2-64aBCx^2.
\]
Give $(z,a,b,c,e,f)$ the weights $(8,16,11,22,18,26)$.
Exact expansion, recorded in a $229$-term certificate, gives
\[
 [N_c^\circ]_{63}=8bc^2z,\qquad
 [N^\circ]_{126}=64b^2c^4(z^2-a).
\]
Substitute
\[
 z=\overline Z-MV-MC_f,\quad a=A,\quad b=MC_f,\quad c=MV,
\]
\[
 e=M^2(I^2-1),\qquad f=-\frac{M^2(V^2-1)}\delta.
\]
These are exact because $C_f^2-D_f=I^2-1$ and
$(MV)^2-C_{\rm num}/\delta=-M^2(V^2-1)/\delta$. Multiplication by
$\delta^2$ and $\delta^4$ gives the actual inner and outer norm.
The power of $f$ is at most four, so the cleared identity is polynomial.
Since $\deg\delta=18$, the exact residual degrees are
\[
 \deg N_c=99,\qquad \deg N_8=198,\qquad\deg A_1=13,
 \qquad\deg P_{11}=211.
\]

For an attained-degree certificate on the actual fixed tuple, put
\begin{align*}
 q_0&=3\beta j,\qquad k_0=\eta+\zeta,\qquad I_0=ik_0^2Y^2,\\
 D_{f,0}&=X^2Y^2I_0^2,\qquad Q_0=X^4Y^4I_0^2,\qquad C_{f,0}=XYI_0,\\
 L_0&=k_0(T+F+\alpha),\qquad\ell_0=Xk_0Y^2,\qquad R_0=q_0^3(q_0-F),\\
 B_{t,0}&=(K_0q_0+X)(q_0-F-\alpha)-q_0F.
\end{align*}
Then
\begin{align*}
 [N_8]_{198}&=64Q_0^4q_0^{32}D_{f,0}^4R_0^4C_{f,0}^2
                 L_0B_{t,0}(L_0B_{t,0}-2\ell_0),\\
 [A_1]_{13}&=32q_0^2k_0X^2Y^4(4q_0^2Y\gamma+R_0).
\end{align*}
The coefficient of
\[
 j^{48}i^{18}\eta^{39}X^{30}Y^{67}T^2F^4\alpha^2\gamma
\]
is $2^{13}(3\beta)^{48}>0$, independently of $K_0,M_C,M_F$.
The literal graph has $64$ multiplications, $31$ additions, and $17$
subtractions. Scalar multiplication and squaring are charged; fixed
numeral leaves and register reuse are free. The two-residual system has
$109$ operations, and the single equation has $112$.

The separate exact audit checks every named polynomial contract, the exact
double-center cancellation, the eight-conjugate identity, radical-extraction
identities, the attained leading coefficient, liveness, and operation
counts. The complete compact44 payload is verified and independently
reconstructed from its finite recipe. Normal and optimized Python runs
pass with explicit checks.

Combining soundness and completeness with the named inherited source
proves the ordinary mathematical consistency equivalence on the stated
positive domain. It does not lower the existing nine-positive variable
count, prove global minimality, or provide a new end-to-end formal
calibration. The new content is the offset, the resulting early sign
argument, its canonical completeness proof, and the degree-$211$ endpoint.



\section{Alternative direct proof sources}
\label{sec:alternative}
\subsection{Optimal natural selector weights}
\label{sec:natural-selectors}

This section gives a completed local improvement to the direct
commutative-certificate architecture.  Its global bounded closure has not
been compiled, so the local dimensions below are not a new final
consistency-equation variable count.

\begin{theorem}[Natural interpolation weights]
Let $k\geq1$, and let $G_0,\ldots,G_{k-1}$ be integer polynomials
nonnegative on natural inputs.  Set
\[
 g_k(r)=\prod_{j=0}^{k-1}(r-j),\qquad
 L_i(r)=\prod_{\substack{0\leq j<k\\j\ne i}}(r-j),
\]
and define
\[
 W_i(r)=
 \begin{cases}
 L_i(r),&k-1-i\text{ even},\\
 (r-k)L_i(r),&k-1-i\text{ odd}.
 \end{cases}
\]
Then $g_k$ and every $W_i$ are nonnegative on $\mathbb N$, and
\[
 \exists r\in\mathbb N:\quad
 g_k(r)+\sum_{i=0}^{k-1}W_i(r)G_i=0
 \quad\Longleftrightarrow\quad
 \bigvee_{i=0}^{k-1}G_i=0.
\]
Each $W_i$ has the smallest possible degree among real polynomials
nonnegative on $\mathbb N$, positive at $i$, and zero at every other
integer in $\{0,\ldots,k-1\}$: degree $k-1$ in the even case and $k$
in the odd case.
\end{theorem}
\begin{proof}
For $r<k$, the factors give $g_k(r)=0$ and $W_i(r)=0$ unless $r=i$.
At $r=i$, the sign of $L_i(i)$ is $(-1)^{k-1-i}$; the optional factor
$i-k$ changes the negative sign.  For $r\geq k$, all relevant factors
are nonnegative, and $g_k(r)>0$.  Thus every summand in the displayed
equation is nonnegative, a zero forces $r<k$, and its unique surviving
branch weight is positive.  This proves both directions.

For optimality, $k-1$ required distinct roots force degree at least
$k-1$.  At that degree a candidate equals $cL_i$; nonnegativity at
large natural arguments forces $c>0$.  Its value at $i$ has the required
sign exactly in the even case.  The odd case therefore needs degree at
least $k$, attained by the stated construction.
\end{proof}

The incoming Round11 source uses twelve nonnegative Horn branch
polynomials $h_0,\ldots,h_{11}$, with respective degrees $2$, ten
copies of $4$, and $8$ for the bundle branch $h_9$.  The truth branch is
$h_0=p^2+q^2+v^2$.  Retain $M=M_0+1$ and the accumulator residuals
\begin{align*}
 E&=S'(1+vM)-S(1+pM)(1+qM),\\
 F&=R'-R\bigl((1+p)(1+q)+v\bigr).
\end{align*}
At a truth step $p=q=v=0$, the equations say $S'=S$ and $R'=R$.
All such steps can therefore be deleted from a finite certificate,
preserving its endpoints and final height.  The existential trace length
and its finite CRT code are replaced by those of the shortened trace.
The distinguished endpoints are $S_0=M+1$ and $S_N=1$, with $M>0$, so a successful trace cannot consist entirely of deleted steps. Conversely every trace without these steps is already an old trace.
This establishes equality of the global existential source predicates;
it does not assert equality of their local zero sets at a fixed step.

Group the remaining eleven factors as
\[
 (G_0,\ldots,G_5)=
 (h_9,h_1h_2,h_3h_4,h_5h_6,h_7h_8,h_{10}h_{11}).
\]
Every group has degree eight.  The resulting literal local polynomial is
\[
 \mathcal L=g_6(r)+\sum_{i=0}^{5}W_i(r)G_i+E^2+F^2.
\]
Its degree is exactly $14$, with $21$ local natural witnesses and five
exposed state/parameter coordinates.  The former three-bit construction
uses $23$ local witnesses at the same degree.  Exact sparse expansion of
the supplied rows gives $50\,461$ monomials; the leading term of
$\mathcal L(t,\ldots,t)$ is $723072t^{14}$.  The source JSON, selector
weights, entire expanded polynomial, and a standard-library verifier are
included.  The exact degree assertion is therefore independently
reproducible.

There is a second improvement using two ternary selectors $u,w$.  The
$k=3$ weights are
\[
 U_0(x)=(x-1)(x-2),\quad
 U_1(x)=x(x-2)(x-3),\quad
 U_2(x)=x(x-1),
\]
of degrees $2,3,2$.  Use the group array
\[
 (G_{ij})=
 \begin{pmatrix}
 h_9&h_5&h_1h_2\\
 h_6&h_7&h_8\\
 h_3h_4&h_{10}&h_{11}
 \end{pmatrix}
\]
and the polynomial
\[
 \mathcal L_{12}=g_3(u)+g_3(w)
 +\sum_{i,j=0}^{2}U_i(u)U_j(w)G_{ij}+E^2+F^2.
\]
The degree-eight groups occupy corners with weight degree four.  Every
other group contributes at degree at most ten.  Thus this source has
$22$ local witnesses and degree exactly $12$; its $40\,990$-monomial
expansion has leading coefficient $627840$ under uniform substitution.
Nonnegativity and the two selector penalties prove the same local and
global source equivalence.  The new local tradeoffs are consequently
$(21,14)$ and $(22,12)$ in place of $(23,14)$.

The incoming index-free inert-prime closure has conditional coordinate
count $2m+9+2f+p+(m+1)b$, for local witness dimension $m$ and chosen
factorial, power, and binomial graph costs $f,p,b$.  Holding the remaining
graph conventions fixed, the change $m=23\longrightarrow21$ saves
$4+2b$ coordinates at that interface.  This is a conditional compiler
accounting identity; it does not replace materialization of the entire
direct consistency polynomial.  The degree-twelve alternative instead
saves $2+b$ coordinates under that same convention.

\subsection{Exact local arithmetic graphs}

A second improvement exploits the natural domain directly.  For natural
$x_1,\ldots,x_s$, their sum vanishes exactly when they all vanish.
Hence a mask $(\sum x_i)F^2$ can replace $(\sum x_i^2)F^2$ without
changing its zeros.  In particular,
\[
 C(n,p)=(n+p)(n-p-1)^2
\]
is a nonnegative cubic with zeros exactly at $p=\operatorname{pred}(n)$.
The corresponding signed quadratic residual $(n+p)(n-p-1)$ has the same
zero set and may be squared when a strict residual-system compiler is
required.

More significantly, the quadratic polynomial
\[
 Z(n,b,p)=(n+b-p-1)^2+bp
\]
jointly represents predecessor and zero test:
\[
 Z=0\quad\Longleftrightarrow\quad
 p=\operatorname{pred}(n),\qquad b=\mathbf1_{n=0}.
\]
Indeed $bp=0$ and $n+b=p+1$.  If $b=0$, then $n=p+1\geq1$; if
$p=0$, then $n+b=1$, giving only $(n,b)=(0,1)$ or $(1,0)$.
The converse is immediate.  Hiding $p$ gives a one-witness zero-test
graph of degree two.  If both outputs are already present, no extra
witness is needed.

For a conditional, one may use the one-witness quartic
\[
 n(c-b)^2+\bigl((c-a)(n-s-1)\bigr)^2,
\]
or the two-witness cubic
\[
 Z(n,t,p)+t(c-a)^2+n(c-b)^2.
\]
Their zeros project exactly to $c=a$ when $n=0$, and $c=b$ otherwise.
For the quartic, $n=0$ makes $n-s-1$ nonzero for every natural $s$;
$n>0$ forces the first branch directly.  For the cubic, the joint graph
fixes the exact branch bit.  Completeness uses the canonical predecessor
and bit.

Keeping the incoming two-witness max-square pairing polynomial $Q$
unchanged, the tree projections use $Q(p,a,b;t,s)+C(n,p)$.  The actual
found-marker list head uses
\[
 Q(p,a,b;t,s)+(n+p+c)(n-p-1)^2+n(c-a-1)^2.
\]
At $n=0$ the middle term forces $p=c=0$; at $n>0$ it forces $p=n-1$
and the last term forces $c=a+1$.  Thus nil and the stored-zero marker
are both preserved.  All naturalness assumptions are explicit.

\begin{center}
\begin{tabular}{lccc}
\hline
Local graph & Hidden witnesses & Degree & Tree operations\\
\hline
Predecessor & $0$ & $4\to3$ & $10\to8$\\
Zero test & $1$ & $4\to2$ & $14\to10$\\
Tree projection/list tail & $4$ & $4$ & $49\to47$\\
Conditional & $1$ & $4$ & $16\to15$\\
List head & $5$ & $4$ & $59\to55$\\
\hline
\end{tabular}
\end{center}

The comparison is against the incoming Joint Guards package and counts
literal binary-expression trees, including repeated square operands and
the same final addition of zero.  Full expansions and separate shared-DAG
counts are supplied.  These are exact local graph improvements, without
a claimed end-to-end compiler or object-level formal certificate.


\section{Logical calibration and what is verified}

All completed equations have literal finite arithmetic descriptions with
fixed integer coefficients. The ten-integer group retains the Round11
H1868 tuple; the integer residual group retains the coherent six-numeral
joint-guard tuple; the positive group retains the complete compact44 tuple.
These tuples must not be mixed. A symbol that appears in one group need
not denote the same integer in another.

The completed reductions prove the following chain:
\[
 \text{new zero exists}\quad\Longleftrightarrow\quad
 \text{calibrated predecessor zero exists}\quad\Longleftrightarrow\quad
 \text{finite ZFC refutation exists}.
\]
For the polynomial shears and monic eliminations, the first equivalence is
an explicit zero-set bijection. The positive-domain component additionally
uses the proved sign-recovery and canonical-witness estimates. The second
equivalence imports the inherited native/compiler/calculus theorem. Its
ordinary ZFC calibration is a stated mathematical theorem, not a conclusion
drawn from an external Lean theorem. The pinned primary K215 proof
\cite{k215} explicitly establishes
\[
 \ZFC\vdash
 \bigl(\operatorname{Halt}(K215,\text{blank})
       \longleftrightarrow\neg\Con(\ZFC)\bigr).
\]
It composes the dominated-calculus refutation theorem, a quantified
source-to-machine simulation, and the exact finite blank startup.
The native arithmetic representation and cited Pell/compiler proofs
then connect this fixed halting predicate with the predecessor equations.
Our integral polynomial transformations and sign argument compose with
those ordinary proofs within ZFC, giving \eqref{eq:goal}.

The distinction from the remaining object-calculus work is explained in
the retained source-63 scope review \cite{zfcscope}. The ordinary argument
codes finite runs and derivations in ambient ZFC and uses uniform
induction over its entire $\omega$. Its model-theoretic steps use set
models, their satisfaction relations, and first-order completeness; they
do not assume transitive models or external well-foundedness. Thus the
claim uses the identified ordinary mathematical calibration as a
dependency. It does not promote correctness on external numerals to
correctness on every internal natural number. No appeal to the actual
truth of $\Con(\ZFC)$ is made at any stage.

There are three separate verification claims:
\begin{enumerate}
\item The new universal algebraic and arithmetic implications are proved
      in this manuscript, with exact coefficient certificates where an
      expansion is too long to print.
\item The supplied programs check literal polynomial transcription,
      coefficient payloads, graph topology, exact algebra, and degree
      witnesses. Finite signed evaluations are secondary regression
      controls; they are not proofs of an unbounded semantic equivalence.
\item A fresh end-to-end Lean/Rocq proof, or an explicit sealed
      object-language ZFC derivation, is not supplied. The repository's
      \file{OBJECT_ZFC.md} correctly distinguishes an external equivalence
      from an internal derivation. This delivery does not close those
      formalization obligations by renaming them.
\end{enumerate}
The direct-proof source in Section~\ref{sec:alternative} has a different
boundary: its local source theorem is proved, but it is not substituted for
the completed native source without constructing and measuring the full
bounded-quantifier elimination. It therefore does not support a claimed
new global number of variables or polynomial degree.

\section{Reproduction and integration}

\file{README.md} gives commands and the file map. Run the root
\file{verify.py} from any working directory. Standard-library checkers are
provided for the ten-integer and positive components. The integer
elimination and independent symbolic audits document their SymPy
requirements explicitly; the alternative generators need only the standard
library. All acceptance checks use explicit exceptions,
rather than depending on Python's removable \file{assert} statements.
The root check records which commands actually ran and whether they passed.

The archive contains the editable manuscript, compiled PDF, complete new
literal equations, all required actual coefficient payloads, source pins,
the relevant predecessor interfaces, independent audits, licenses, and
notices. The optional article build uses pdfLaTeX and the packages named in
this file. It uses no network connection. The source snapshot and incoming
archive identities are fixed in \file{PROVENANCE.json}.

For repository intake, place the archive's wrapper directory in a new
proposal subdirectory under
\file{Computability/DiophantineConsistency/con-zfc-diophantine-forms}.
No directory number is reserved here. The delivery makes no edit to the
repository's main branch or existing packages. Comparisons should retain
both domains and all measured tradeoffs, particularly the operation cost
of the six-term ten-integer variant.

\section{Further research questions}

\begin{enumerate}
\item \textbf{Go beyond translations of the last coordinate.}
The odd-degree argument proves that $1\,081\,125$ is the floor for
polynomial translations of $m$ with the present four norm inputs fixed.
Can a joint triangular change of $y$, $z$, or a source parameter lower
$\dd(\Sigma^4a_0a_1)$ without increasing the circuit cost excessively?
This question has a clear obstruction and a concrete target.

\item \textbf{Reduce the cost of the six-term shear.}
Its factored expression adds $41$ live operations. Can the $26$-monomial
polynomial share additional powers and products with $D,F,G,H$ and the
source prefix? A claimed improvement should be accompanied by a literal
graph and exact identity, rather than a heuristic cost estimate.

\item \textbf{Audit other radical norms for removable leading squares.}
The method applies whenever a large radicand has an integral polynomial
part. For unrestricted coordinates, completeness is automatic under a
triangular shear. Positive coordinates require a new range argument; a
formal square-root expansion by itself cannot supply that argument.

\item \textbf{Improve the degree-211 positive endpoint.}
The new transport offset supplies a strict sign before native code
recovery. Can a smaller positive offset do the same uniformly? Can the
majorant or root center be simplified while preserving the order of the
sign-recovery argument? A negative or zero intermediate must not be
excluded by an unstated positivity assumption.

\item \textbf{Continue the monic elimination search with degree control.}
The $18$-integer form has six residuals and degree $88$. Search for other
polynomial changes that expose a monic row while keeping the largest
residual below degree $44$. A Gröbner-style symbolic search can assist,
but the final certificate should give explicit integral restoration maps.

\item \textbf{Compile the direct certificate source and measure it.}
The natural-domain selector and functional graph improvements reduce
local cost without requiring K215. The next necessary object is a complete
ordinary-polynomial constructor with its actual witness, gate, and
coefficient counts, including bounded-universal elimination. Local counts
alone are not a size comparison with the finished equations.

\item \textbf{Close the internal-ZFC proof boundary.}
Formalize the actual coefficient-bound native compiler, its source
calibration, and the new simple shear/monic-elimination lemmas in the
repository's sealed object-language framework. The new local algebra is
short; uniform correctness of the inherited coding over the internal
natural numbers of every ZFC model remains the substantial work.
\end{enumerate}

\begin{thebibliography}{99}
\raggedright
\bibitem{jones1982}
J. P. Jones, \emph{Universal Diophantine equation}, Journal of Symbolic
Logic 47 (1982), 549--571.
\href{https://doi.org/10.2307/2273588}{doi:10.2307/2273588}.

\bibitem{sun2026}
Z.-W. Sun, \emph{Ten unknowns for Hilbert's tenth problem over the integers},
arXiv:2609.01594v2, 2026.
\url{https://arxiv.org/abs/2609.01594}.

\bibitem{repo}
V. Reshetnikov, \emph{ProveIt}, inspected repository snapshot
\file{c1a102936aca437e9999530a0111e6ddab56dc6d}.
\url{https://github.com/VladimirReshetnikov/ProveIt}.
See the DiophantineConsistency catalogue and its source-pinned research
packages; their individual notices are retained with reused artifacts.

\bibitem{k215}
ProveIt contributors, \emph{Ordinary proof of the combined215 machine},
\file{Computability/BusyBeaver/Research/ZFC/candidate-k215/proof.md},
at the inspected repository snapshot. Retained verbatim as
\file{calibration/K215_CALIBRATION_PROOF.md}; its Git blob is
\file{5b62cecff23a798a4b34255d1f843e72ad11b623}.

\bibitem{zfcscope}
ProveIt source 63, \emph{Ordinary ZFC provability and the separate
object-calculus obligation},
\file{63-sign-recovery/reviews/ZFC_INTERNALIZATION_SCOPE.md}.
Retained verbatim under \file{calibration/}, with its original source
snapshot and scope qualifications.

\bibitem{round11}
ProveIt incoming package,
\emph{Structural reductions for Diophantine consistency equations,
Round eleven}, archive
\file{ConZFC_Structural_Reductions_Round11_2026-10-24.zip},
especially \file{lucas_h1868/PROOF.md}. Exact archive and member pins are
recorded in the accompanying provenance files.

\bibitem{joint}
ProveIt incoming package,
\emph{Joint sign guards and Pell descent for smaller consistency equations},
\file{ZFC_Diophantine_Joint_Guards_2026-10-24.zip}.

\bibitem{centering}
ProveIt incoming package,
\emph{Joint guards and two Pell-root translations},
\file{con-zfc-budget-and-centering-2026-10-23.zip}.

\bibitem{radical}
ProveIt incoming package,
\emph{Radical majorants, guard compression, and Pell descent},
\file{con-zfc-radical-guard-descent-2026-10-23.zip}.

\bibitem{source68}
ProveIt source 68, \emph{Joint positivity and free Lucas coordinates},
\file{68-joint-positivity}, notably its direct certificate source.

\bibitem{source71}
ProveIt source 71, \emph{Calibrated determinant guards},
\file{71-determinant-guards}.

\bibitem{source75}
ProveIt source 75, \emph{Signed-power recovery and inert-prime certificates
for ZFC consistency}, \file{75-signed-bootstrap}.
\end{thebibliography}
\end{document}
